% Generated draft: requires mathematical and source review.
\documentclass[11pt]{article}
\usepackage[margin=25mm]{geometry}
\usepackage{fontspec}
\setmainfont{Latin Modern Roman}
\usepackage{amsmath,amssymb,mathtools,unicode-math}
\setmathfont{Latin Modern Math}
\usepackage{enumitem,needspace,xurl,hyperref}
\setlength{\parindent}{0pt}
\setlength{\parskip}{5pt}
\setlength{\emergencystretch}{3em}
\newcommand{\N}{\mathbb{N}}
\newcommand{\Z}{\mathbb{Z}}
\newcommand{\Q}{\mathbb{Q}}
\newcommand{\R}{\mathbb{R}}
\newcommand{\C}{\mathbb{C}}
\newcommand{\F}{\mathbb{F}}
\newcommand{\A}{\mathbb{A}}
\newcommand{\PP}{\mathbb{P}}
\newcommand{\E}{\mathbb{E}}
\newcommand{\eps}{\varepsilon}
\newcommand{\dd}{\,\mathrm{d}}
\begin{document}
\section*{UnsolvedMath ID: 20003273}
\textbf{Source problem code:} catalogue entry
\subsection{Definitions and mathematical statement}
D. Fisher's Question 6 asks whether measure-rigidity results for higher-rank abelian groups yield new obstructions to smooth or continuous actions of such groups on compact manifolds. For example, for $k\ge2$ one may consider actions $\alpha:\Z^k\to\operatorname{Diff}^\infty(M)$ or continuous actions of $\Z^k$ or $\R^k$ on compact $M$. An action is effective if only the identity group element acts identically. The source does not specify effectiveness, a genuine-rank condition, invariant measures, or positive entropy, so its target is a research direction rather than a single universal nonexistence statement.

The source notes existing obstructions for smooth volume-preserving actions due particularly to H. Hu and A. Katok. It explains why arithmetic-quotient arguments for semisimple groups and lattices do not transfer directly: those arguments use both Ratner's theorem and superrigidity for merely measurable cocycles. The paragraph invokes Dye's theorem [43] and generalizations [154, 37] as an amenable orbit-equivalence obstacle to the needed measurable-cocycle rigidity. This is the source's contextual limitation, not an assertion excluding every possible restricted cocycle theorem.

The attached aid reports affirmative topological restrictions, including exclusion of spheres, for maximal-rank actions under strong positive-entropy and nonuniform-rigidity hypotheses, attributed to Katok and Rodriguez Hertz. Those hypotheses and full citations are not supplied, so this cannot be stated as a restriction on all higher-rank abelian actions.

In the opposite direction, its reported universality statement is: for every positive-dimensional closed connected smooth manifold $M$ and every $k\ge2$, there exist effective smooth $\R^k$- and $\Z^k$-actions for which every element has zero entropy. The $\R^k$ orbits have dimension at most one, with the discrete examples given by restriction; for $\dim M\ge2$, the actions can preserve any prescribed smooth volume. Thus these examples do not supply genuine higher-dimensional orbit dynamics, but they rule out an obstruction to effective actions based on the underlying manifold alone.

The aid also reports this entropy gate. Let $\theta>0$, let $\alpha:\Z^k\to\operatorname{Diff}^{1+\theta}(M)$ preserve a smooth probability volume $\mu$, and let $f=\alpha(a)$ have a $Df$-invariant rank-$r$ bundle $E$, $r\ge1$, with
\[
 \|Df^j(v)\|\ge C\lambda^j\|v\|
 \quad(j\ge1,\ x\in M,\ v\in E_x),
 \qquad C>0,\quad\lambda>1.
\]
The reported Pesin-formula transfer gives
\[
 h_\mu(f)\ge r\log\lambda>0.
\]
In the full-action ergodic decomposition $\mu=\int\mu_\omega\,\dd\eta(\omega)$, entropy affinity yields positive-entropy $f$-components on a set of positive $\eta$-measure. A measure-rigidity theorem can yield a topological obstruction only when its remaining rank, regularity, and dynamical hypotheses apply to the relevant action and measures.

The appended heading ``3. Measure rigidity'' and section 3.1, about $G/\Gamma$ and a subgroup generated by one-parameter unipotent subgroups, begin the next problem's context. They are not hypotheses of Fisher's question. All stated progress is from the supplied dated machine-generated aid and requires independent verification.
\subsection{Short English statement}
Can higher-rank abelian measure rigidity produce new restrictions on manifold actions? The supplied progress requires positive entropy and genuine rank: effective zero-entropy actions with at most one-dimensional orbits exist on every closed connected manifold of positive dimension, so no unconditional manifold obstruction follows.
\subsection{Sources}
\begin{itemize}
\item D. Fisher, Question 6, American Institute of Mathematics, \emph{Emerging applications of measure rigidity}, page 4 according to the supplied verification, \url{https://aimath.org/WWN/measrigid/measrigid.pdf}. Source names H. Hu, A. Katok and Ratner; it cites Dye [43] and generalizations [154, 37]. The spill into section 3.1 cites [159] and exposition [144] concerning unipotent measure rigidity. Full bibliography entries are not supplied.
\item The research summary names Katok--Rodriguez Hertz restrictions for maximal-rank actions and invokes Pesin's entropy formula, without supplying complete bibliographic references or the full maximal-rank theorem hypotheses.
\item Local UnsolvedMath registry, permanent ID $20003273$, source code \texttt{catalogue entry}, canonical location \texttt{aim-other-notes.json}, \texttt{notes[79]}, dated 14 August 2026. Research attempt 1: partial result, proved rigor, low novelty confidence; machine-generated and not peer reviewed.
\end{itemize}
\end{document}
